% Lemmas on Feasible Error Region for Gradient-Constrained Compression
% Include this file in your main document or compile standalone

\documentclass{article}
\usepackage{amsmath,amssymb,amsthm}

\newtheorem{lemma}{Lemma}
\newtheorem{definition}{Definition}

\begin{document}

\begin{definition}[Linear Operators on Discrete Fields]
\label{lem:linear-operator}
Let $\mathbf{x} \in \mathbb{R}^n$ be a data array defined on a $d$-dimensional regular grid with dimensions $n_1 \times n_2 \times \cdots \times n_d$, where $n = \prod_{j=1}^d n_j$. The data values are stored in a flattened vector using row-major ordering.

Any linear operation on $\mathbf{x}$ that produces an output $\mathbf{y} \in \mathbb{R}^n$ of the same dimension can be represented as
\begin{equation}
    \mathbf{y} = \mathbf{A}\mathbf{x}
\end{equation}
for some matrix $\mathbf{A} \in \mathbb{R}^{n \times n}$. The choice of boundary condition (circular, symmetric, replicate, or zero) affects the values of the matrix entries but does not alter the linearity of the operation.
\end{definition}

\begin{definition}[Multi-Operator Constraints]
\label{def:multi-operator-constraints}
Given $m$ linear operators $\mathbf{A}_1, \ldots, \mathbf{A}_m \in \mathbb{R}^{n \times n}$ with corresponding per-element bound vectors $\mathbf{b}_1, \ldots, \mathbf{b}_m \in \mathbb{R}^n_{>0}$, and a box constraint bound $\mathbf{d} \in \mathbb{R}^n_{>0}$, the feasible region is:
\begin{equation}
    \mathcal{S} = \left\{ \mathbf{x} \in \mathbb{R}^n \;\middle|\; |\mathbf{x}| \le \mathbf{d} \text{ and } |\mathbf{A}_k \mathbf{x}| \le \mathbf{b}_k \text{ for all } k \in \{1, \ldots, m\} \right\}
\end{equation}
\end{definition}

\begin{lemma}[Convexity of the Feasible Region]
\label{lem:feasible-region}
The feasible region $\mathcal{S}$ defined in Definition~\ref{def:multi-operator-constraints} is convex.
\end{lemma}

\begin{proof}
Let $\mathbf{x}_1, \mathbf{x}_2 \in \mathcal{S}$ and consider any convex combination $\mathbf{x}_t = t\mathbf{x}_1 + (1-t)\mathbf{x}_2$ for $t \in [0, 1]$.

\textbf{Box constraint:} Since $|\mathbf{x}_1| \le \mathbf{d}$ and $|\mathbf{x}_2| \le \mathbf{d}$, by the triangle inequality:
\begin{equation}
    |\mathbf{x}_t| = |t\mathbf{x}_1 + (1-t)\mathbf{x}_2| \le t|\mathbf{x}_1| + (1-t)|\mathbf{x}_2| \le t\mathbf{d} + (1-t)\mathbf{d} = \mathbf{d}.
\end{equation}

\textbf{Operator constraints:} For each $\mathbf{A}_k$, by linearity:
\begin{equation}
    |\mathbf{A}_k \mathbf{x}_t| = |t\mathbf{A}_k \mathbf{x}_1 + (1-t)\mathbf{A}_k \mathbf{x}_2| \le t|\mathbf{A}_k \mathbf{x}_1| + (1-t)|\mathbf{A}_k \mathbf{x}_2| \le t\mathbf{b}_k + (1-t)\mathbf{b}_k = \mathbf{b}_k.
\end{equation}

Thus $\mathbf{x}_t \in \mathcal{S}$, proving convexity.
\end{proof}

\subsection*{Objection}

Let $\boldsymbol{\varepsilon}_0 = \tilde{\mathbf{x}} - \mathbf{x}$ denote the error vector after compressing and decompressing $\mathbf{x}$ by an error-bounded lossy compressor, where $\tilde{\mathbf{x}}$ is the reconstructed data. The compressor guarantees that the initial error $\boldsymbol{\varepsilon}_0$ satisfies the box constraint $|\boldsymbol{\varepsilon}_0| \le \mathbf{d}$, but may violate the operator constraints $|\mathbf{A}_k \boldsymbol{\varepsilon}_0| \le \mathbf{b}_k$. Our objective is to find a corrected error $\hat{\boldsymbol{\varepsilon}} \in \mathcal{S}$ that lies in the feasible region (Definition~\ref{def:multi-operator-constraints}) while minimizing the modification to $\boldsymbol{\varepsilon}_0$.

\begin{equation}
\begin{aligned}
    \min_{\hat{\boldsymbol{\varepsilon}}} \quad & \|\hat{\boldsymbol{\varepsilon}} - \boldsymbol{\varepsilon}_0\|_1 \\
    \text{subject to} \quad & |\hat{\boldsymbol{\varepsilon}}| \le \mathbf{d}, \\
    & |\mathbf{A}_k \hat{\boldsymbol{\varepsilon}}| \le \mathbf{b}_k \quad \text{for all } k \in \{1, \ldots, m\}.
\end{aligned}
\end{equation}
This formulation is a linear program (LP) that can be solved efficiently. The L1 norm is chosen because it is the convex envelope of the L0 pseudo-norm on bounded domains, making it the tightest convex relaxation for minimizing the number of modified entries. While L1 minimization does not guarantee sparse solutions in general, it empirically tends to produce sparser modifications than L2 minimization, which distributes changes smoothly across all components.

\subsection*{PDHG Solver}

To apply the Primal-Dual Hybrid Gradient (PDHG) algorithm, we first transform the problem into standard form.

\textbf{Step 1: Variable Substitution.} Let $\boldsymbol{\delta} = \hat{\boldsymbol{\varepsilon}} - \boldsymbol{\varepsilon}_0$ denote the modification vector. The problem becomes:
\begin{equation}
\begin{aligned}
    \min_{\boldsymbol{\delta}} \quad & \|\boldsymbol{\delta}\|_1 \\
    \text{subject to} \quad & |\boldsymbol{\delta} + \boldsymbol{\varepsilon}_0| \le \mathbf{d}, \\
    & |\mathbf{A}_k \boldsymbol{\delta} + \mathbf{A}_k \boldsymbol{\varepsilon}_0| \le \mathbf{b}_k \quad \text{for all } k \in \{1, \ldots, m\}.
\end{aligned}
\end{equation}

\textbf{Step 2: Stack Operators.} Define the stacked operator matrix and offset vector:
\begin{equation}
    \mathbf{A} = \begin{bmatrix} \mathbf{A}_1 \\ \vdots \\ \mathbf{A}_m \end{bmatrix} \in \mathbb{R}^{mn \times n}, \quad
    \mathbf{c} = \begin{bmatrix} \mathbf{A}_1 \boldsymbol{\varepsilon}_0 \\ \vdots \\ \mathbf{A}_m \boldsymbol{\varepsilon}_0 \end{bmatrix} = \mathbf{A} \boldsymbol{\varepsilon}_0 \in \mathbb{R}^{mn}, \quad
    \mathbf{b} = \begin{bmatrix} \mathbf{b}_1 \\ \vdots \\ \mathbf{b}_m \end{bmatrix} \in \mathbb{R}^{mn}
\end{equation}
The operator constraints become a single constraint $|\mathbf{A}\boldsymbol{\delta} + \mathbf{c}| \le \mathbf{b}$.

\textbf{Step 3: Linearize Absolute Value Constraints.} The constraint $|\mathbf{A}\boldsymbol{\delta} + \mathbf{c}| \le \mathbf{b}$ expands to:
\begin{equation}
    \mathbf{A}\boldsymbol{\delta} \le \mathbf{b} - \mathbf{c} \quad \text{and} \quad -\mathbf{A}\boldsymbol{\delta} \le \mathbf{b} + \mathbf{c}
\end{equation}
Similarly, the box constraint $|\boldsymbol{\delta} + \boldsymbol{\varepsilon}_0| \le \mathbf{d}$ becomes:
\begin{equation}
    \boldsymbol{\delta} \le \mathbf{d} - \boldsymbol{\varepsilon}_0 \quad \text{and} \quad -\boldsymbol{\delta} \le \mathbf{d} + \boldsymbol{\varepsilon}_0
\end{equation}
Define the asymmetric box bounds $\mathbf{l} = -\mathbf{d} - \boldsymbol{\varepsilon}_0$ and $\mathbf{u} = \mathbf{d} - \boldsymbol{\varepsilon}_0$, so the box constraint is $\mathbf{l} \le \boldsymbol{\delta} \le \mathbf{u}$.

\textbf{Step 4: LP Reformulation.} Split the variable $\boldsymbol{\delta} = \boldsymbol{\delta}^+ - \boldsymbol{\delta}^-$ where $\boldsymbol{\delta}^+, \boldsymbol{\delta}^- \ge 0$. Introduce slack variables $\mathbf{s}_1, \mathbf{s}_2 \ge 0$ for the operator constraints. The LP becomes:
\begin{equation}
    \min_{\tilde{\boldsymbol{\delta}}} \mathbf{c}_{\text{obj}}^T \tilde{\boldsymbol{\delta}} \quad \text{subject to} \quad \tilde{\mathbf{A}}\tilde{\boldsymbol{\delta}} = \tilde{\mathbf{b}}, \quad \mathbf{0} \le \tilde{\boldsymbol{\delta}} \le \mathbf{U}
\end{equation}
where the block matrices are:
\begin{equation}
    \tilde{\boldsymbol{\delta}} = \begin{bmatrix} \boldsymbol{\delta}^+ \\ \boldsymbol{\delta}^- \\ \mathbf{s}_1 \\ \mathbf{s}_2 \end{bmatrix}, \quad
    \mathbf{c}_{\text{obj}} = \begin{bmatrix} \mathbf{1}_n \\ \mathbf{1}_n \\ \mathbf{0}_{mn} \\ \mathbf{0}_{mn} \end{bmatrix}, \quad
    \mathbf{U} = \begin{bmatrix} \mathbf{u} \\ -\mathbf{l} \\ \infty \\ \infty \end{bmatrix}
\end{equation}
\begin{equation}
    \tilde{\mathbf{A}} = \begin{bmatrix} \mathbf{A} & -\mathbf{A} & \mathbf{I}_{mn} & \mathbf{0} \\ -\mathbf{A} & \mathbf{A} & \mathbf{0} & \mathbf{I}_{mn} \end{bmatrix}, \quad
    \tilde{\mathbf{b}} = \begin{bmatrix} \mathbf{b} - \mathbf{c} \\ \mathbf{b} + \mathbf{c} \end{bmatrix}
\end{equation}

\textbf{Step 5: PDHG Algorithm.} Construct the Lagrangian with dual variables $\mathbf{y} \in \mathbb{R}^{2mn}$:
\begin{equation}
    L(\tilde{\boldsymbol{\delta}}, \mathbf{y}) = \mathbf{c}_{\text{obj}}^T \tilde{\boldsymbol{\delta}} - \mathbf{y}^T(\tilde{\mathbf{A}}\tilde{\boldsymbol{\delta}} - \tilde{\mathbf{b}})
\end{equation}
For iteration $k \ge 0$ with step sizes $\tau, \sigma > 0$:
\begin{align}
    \text{Primal update:} \quad & \tilde{\boldsymbol{\delta}}_{k+1} = \min\left(\mathbf{U}, \max\left(\mathbf{0}, \tilde{\boldsymbol{\delta}}_k - \tau (\mathbf{c}_{\text{obj}} - \tilde{\mathbf{A}}^T \mathbf{y}_k)\right)\right) \\
    \text{Extrapolation:} \quad & \bar{\boldsymbol{\delta}}_{k+1} = 2\tilde{\boldsymbol{\delta}}_{k+1} - \tilde{\boldsymbol{\delta}}_k \\
    \text{Dual update:} \quad & \mathbf{y}_{k+1} = \mathbf{y}_k + \sigma (\tilde{\mathbf{A}}\bar{\boldsymbol{\delta}}_{k+1} - \tilde{\mathbf{b}})
\end{align}

\textbf{Step 6: Recovery.} Upon convergence, recover the original solution:
\begin{equation}
    \boldsymbol{\delta} = \boldsymbol{\delta}^+ - \boldsymbol{\delta}^-, \quad \hat{\boldsymbol{\varepsilon}} = \boldsymbol{\delta} + \boldsymbol{\varepsilon}_0
\end{equation}

\end{document}
